\section{Билет 42. Закон равнораспределения энергии по степеням свободы. Средняя кинетическая энергия поступательного движения одноатомной молекулы. Внутренняя энергия газа. Способы изменения внутренней энергии.}

\begin{definition}
Согласно \textbf{закону равнораспределения энергии по степеням свободы}, при термодинамическом равновесии на каждую степень свободы молекулы в среднем приходится энергия $kT/2$. Для одноатомной молекулы ($i=3$) средняя кинетическая энергия поступательного движения равна:
\begin{equation}
    \langle E_{\text{пост}} \rangle = \frac{3}{2} k T. \label{eq:mono_energy_42}
\end{equation}
Абсолютная температура $T$ прямо пропорциональна средней кинетической энергии теплового движения частиц. Термодинамическая температура обращается в нуль, когда средняя кинетическая энергия молекул становится равной нулю.
\end{definition}

\begin{definition}
\textbf{Внутренней энергией} $U$ термодинамической системы называется сумма кинетических энергий хаотического теплового движения всех её микрочастиц и потенциальных энергий их взаимодействия друг с другом.
Для идеального газа потенциальной энергией взаимодействия молекул пренебрегают. Следовательно, внутренняя энергия идеального газа равна сумме энергий всех его молекул:
\begin{equation}
    U = N \langle \epsilon \rangle,
\end{equation}
где $N$ --- общее число молекул, $\langle \epsilon \rangle = \frac{i}{2} k T$ --- средняя энергия одной молекулы.
\end{definition}

\begin{theorem}[Формула для внутренней энергии]
Подставляя $N = \nu N_A$ и используя связь $R = k N_A$, получаем макроскопическое выражение для внутренней энергии произвольной массы идеального газа:
\begin{equation}
    U = \nu N_A \frac{i}{2} k T = \nu \frac{i}{2} R T = \frac{m}{M} \frac{i}{2} R T. \label{eq:internal_energy_42}
\end{equation}
Используя уравнение состояния $pV = \nu RT$, внутреннюю энергию можно выразить через макроскопические параметры:
\begin{equation}
    U = \frac{i}{2} p V. \label{eq:U_pV_42}
\end{equation}
Внутренняя энергия является \textbf{функцией состояния}: её изменение $\Delta U$ зависит только от начального и конечного состояний системы и не зависит от пути перехода.
\end{theorem}

\begin{property}[Способы изменения внутренней энергии]
Существует два принципиально различных способа изменения внутренней энергии термодинамической системы:
\begin{enumerate}
    \item \textbf{Выполнение механической работы} $A$ над системой или самой системой. Механическая работа связана с упорядоченным перемещением макроскопических частей системы. При сжатии газа работа внешних сил положительна ($A_{\text{внеш}} > 0$), внутренняя энергия возрастает. При расширении газ совершает работу над внешними телами ($A > 0$), и его внутренняя энергия убывает.
    \item \textbf{Теплопередача} (теплообмен) --- процесс передачи энергии от одной части системы к другой или от одной системы к другой за счёт хаотического движения молекул (теплопроводность, конвекция, излучение). Количество переданной энергии обозначается $Q$.
\end{enumerate}
Эти два способа количественно связаны \textbf{первым началом термодинамики}:
\begin{equation}
    Q = \Delta U + A. \label{eq:first_law_42}
\end{equation}
Теплота $Q$, сообщаемая системе, расходуется на изменение её внутренней энергии $\Delta U$ и на совершение системой работы $A$ над внешними телами.
\end{property}
